\chapter{Handoff instructions for the next researcher}\label{ch:next}
\section{The task is a new uniform arithmetic bound}
Continue toward either a proof or a counterexample to \eqref{eq:EP}. Do not
spend the next pass merely repackaging the same gap or reporting that the
examples have small fibers. The best-supported starting points are the exact
fixed-cap reduction, primitive-divisor exponent lists, full valuation equations,
and genuine-exchange decomposition. The complete fixed-two-prime proof removes
all two-base exchanges, but leaves the interaction of larger exchanges open.

Before leaning heavily on the newest two-prime theorem, independently audit its
reduced-denominator formula, the bounds $d>f\ge2$ and $A<2d,B<2f$, the
prime-adic inequalities, the integer alignment $\alpha f=\beta d$, and the final
monotonicity identity. Finite regression cannot replace this proof audit.
Likewise, audit the strongest general chain-packing proof if its constant is
used in a new argument. This prioritizes central dependencies rather than
rechecking every already superseded numerical estimate.

\section{Three remaining routes, stated honestly}
\paragraph{Actual exchange packing.}
Prove a uniform bound on the independent-set count of exact minimal exchanges,
without first enumerating the unknown fiber. Retain overlap information when
useful; the product bound in Chapter~\ref{ch:exchanges} can be much stronger
than necessary. Every minimal exchange now has at least three changed bases,
but a fixed lower bound on its support size is insufficient. Look for a proved
relation among the number of independent exchanges, their input-prime scales,
and their total target cost.

\paragraph{Integer reconstruction.}
Use exact row-support tests, integer-valid rounding, certificates with a
controlled positive budget, or a deterministic branching decoder. A successful
proof must bound the worst-case logarithmic choice cost on all targets at every
fixed cap. It need not give a fast algorithm: forced computation and branching
ambiguity are different quantities. The singleton fractional example should
prevent a return to dimension-only optimism.

\paragraph{Arithmetic restrictions beyond potential feedback.}
The signed-graph criterion remains valid after sound target-specific domain
reductions. Counting every unpruned universal positive cycle is conditionally
obstructed by the polynomial-prime family, so do not assume that route's
subpolynomial hypothesis. A useful refinement must quantify which cycles survive
all required valuation capacities and which internally possible fillings fail
on outside-prime rows. The arbitrary fixed-two-prime injectivity theorem is one
example of such arithmetic elimination.

None of these paragraphs is a proved finishing lemma. Their purpose is to make
the required new information explicit.

\section{Counterexample program}
Search for pairwise input-support-disjoint genuine collisions at bounded exponent
cap. If $h(a_i)=h(b_i)=M_i$ and the supports of $a_ib_i$ are pairwise disjoint,
then $2^r$ inputs share target $\prod_{i=1}^rM_i$. A construction with total
logarithmic target cost $O(r\log r)$ would disprove the primary question.
Repeated coprime scaling of one collision does not give this independence.
A large graph cycle without an exact equality also does not count.

For a conjectured intermediate statement, search deliberately for its smallest
counterexample or an infinite obstruction family before polishing a proof around
it. Use the main witnesses in Chapter~\ref{ch:examples}, including the
polynomial-pair targets with many unpruned cycles and unique cap-five inputs.

\section{Computational protocol}
Build complete local domains from the target factorization or the exact
multiplicative-order lists. Include every output-prime row. Choose at most one
local prime-power block per input base, require exact exhaustion of all target
valuations, and prove each pruning operation necessary. Label interrupted
searches incomplete. Use integer or rational arithmetic to certify discovered
witnesses and row weights. Numerical optimization can discover a candidate but
its infeasibility or optimality status is not automatically a proof.

Separate the input-size cutoff, target-size cutoff, allowed prime-factor range,
exponent cap, and divisor-sum roughness restriction in every experiment.
When practical, verify a witness by constructing and summing its divisors, a
method independent of the geometric-sum implementation. For a specified factored
target, exhaustive divisor enumeration is complete because every preimage
divides the target. Do not claim minimality of a target from a particular-fiber
check.

If actual parallel-worker tools are available, assign distinct roles to proof,
source verification, counterexample search, and adversarial audit. Otherwise
perform clearly separated passes. Do not report independent agents or formal
verification that did not occur.

\section{Copy-paste continuation brief}
\begin{statusbox}{Instruction to the receiving agent}
Investigate Erd\H{o}s 1060 using this dossier and its accompanying archive.
The objective is the uniform bound
$\log\max(1,f(N))=o(\log N/\log\log N)$ for $f(N)=\#\{k:k\sigma(k)=N\}$.
No complete proof currently appears in the record.

Start with the fixed-cap reduction, the complete fixed-two-prime rigidity
proof, primitive-divisor exponent lists, and genuine exchanges. Independently
audit any working proof on which your new step depends. Use the strongest
general bound only as a partial result, not as evidence that its constant can
be made arbitrarily small.

Do not assume the unproved low-cost feedback, low-nullity, cheap reconstruction,
or exchange-packing bounds. The dossier includes exact counterexamples and a
conditional obstruction to several versions of those assertions. Preserve
all output-prime valuation equations and distinguish integer solutions from
fractional ones.

Develop a concrete argument far enough to prove a new lemma, find a
counterexample, or identify one precise implication still missing. Continue
revising failed routes rather than returning another generic list of ideas.
Check every quantifier and parameter dependence before claiming a solution.
A proof at every fixed exponent cap would complete the unrestricted problem;
a proof at just one cap would not.

Use reproducible exact tests to challenge conjectured lemmas. State any
unproved assertion explicitly. Return a full proof only when the entire
implication to the uniform little-$o$ estimate is established. Otherwise return
the strongest actual result and its exact remaining gap, with no claim of
completion or proximity.
\end{statusbox}

\section{What a final proof audit must check}
The normalized target coefficient must approach zero for every sufficiently
large $N$. Every residual target in a decomposition must be handled uniformly,
including bounded ones. Only unitary blocks may be automatically canceled.
The exponent-index conventions and primitive-divisor exceptions must match.
No density-one theorem, finite range, conditional prime-pattern prediction,
positive-constant estimate, or real-relaxation fact may be substituted for a
uniform integer-fiber theorem. Newness and correctness are separate questions:
verify the latter before making a publication or priority claim.
